% problem-000334 6 卤化氢加成动力学
\section*{6. 卤化氢加成动力学}
\textit{下列为分问。}

\subsection*{6-1}
对于A为�级，B为�级的反应，在较短时间间隔内 $\Delta t$ ，有：

$
c (\mathrm{AB}) / c (\mathrm{A}) = k c ^ {m - 1} (\mathrm{A}) c ^ {n} (\mathrm{B}) \Delta t
$

实验发现，上述加成反应的�(AB)/�(B)与�(B)⽆关，在保持�(B)不变的条件下，分别取 $\langle p _ { \mathrm { A } }$ 为 ${ \mathsf { I } } { \mathsf { p } } ^ { \ominus }$ 、 $6 p ^ { \ominus }$ $3 p ^ { \ominus }$ $( p ^ { \ominus } = 1 0 0 \mathrm { { k P a } ) \mathrm { { f } } }$ 时，测得 $c ( \mathrm { A B } ) / c ( \mathrm { A } )$ 值之⽐为 $9 { : } 4 { : } 1 _ { \circ }$ 求该反应各反应物的级数和反应总级数（按理想⽓体处理）。

\subsection*{6-2}
以d $c ( \mathrm { A B } ) / \mathrm { d }$ �表⽰该反应速率，写出速率⽅程。

\subsection*{6-3}
设 $\mathrm { \bar { \it { c } } } ( \mathrm { A } )$ 为� mol L−1且保持不变， $c _ { 0 } ( \mathrm { B } ) = 1 \mathrm { m o l } \mathrm { L } ^ { - 1 }$ $c _ { 0 } ( \mathrm { A B } ) = 0$ ，写出 $_ { \mathrm { { \tiny ~ \mathscr { C } } } } ( { \mathrm { { B } } } ) = 0 . 2 5 { \mathrm { m o l } } { \mathrm { { L } } } ^ { - }$ 1时反应所需时间的表达式。

\subsection*{6-4}
有⼈提出该反应的反应历程为：

$
\begin{array}{c} 2 \mathrm{HCl} \xleftarrow [ k _ {- 1} ]{k _ {1}} (\mathrm{HCl}) _ {2} \\ \mathrm{HCl} + \mathrm{C} _ {3} \mathrm{H} _ {6} \xleftarrow [ k _ {- 2} ]{k _ {2}} \mathrm{C} ^ {*} \\ \mathrm{C} ^ {*} + (\mathrm{HCl}) _ {2} \xrightarrow {k _ {3}} \mathrm{CH} _ {3} \mathrm{CHClCH} _ {3} + 2 \mathrm{HCl} \end{array}
$

请根据有关假设推导出该反应的速率⽅程，写出表观速率系数的表达式，将速率⽅程与6-2的结论进⾏⽐较，说明什么问题？

\subsection*{6-5}
实验测得该反应在 $7 0 ~ ^ { \circ } \mathrm { C }$ 时的表观速率系数为 $1 9 ~ ^ { \circ } \mathrm { C }$ 时的 $1 / 3$ ，试求算该反应的表观活化能 $E _ { \mathrm { a ^ { c } } }$ 。
